\subsection{Second proof}

\begin{proof}[Second proof of Theorem \ref{thm:eig-loc-U}]
We keep the notation
\[
f(|z|):=|z|^{2k}e^{-\pi|z|^2},
\qquad
c:=c(\lambda,k),
\]
from the first proof. The argument below is in the spirit of the holomorphic
moving-plane/quadrature-domain methods of \cite{Aharonov-Schiffer-Zalcman,Gustafsson,Karp},
in that we exploit the existence of a holomorphic function on \(U\) with prescribed real
boundary values to recover full radial symmetry. Starting again from \eqref{eq:fourier-delta} and
applying Lemma~\ref{lemma:fourier-to-fboundary}, we obtain a compactly supported distributional
solution of
\begin{equation}\label{eq:potato-second-proof}
\Delta u=f(|z|)\mathbf{1}_U-c\,\delta_0
\qquad\text{in }\mathcal D'(\R^2).
\end{equation}
Using the planar fundamental solution \((2\pi)^{-1}\log|z|\), we may write
\begin{equation}\label{eq:explicit-u-second-proof}
u(z):=\frac{1}{2\pi}\int_U \log|z-z'|\,f(|z'|)\,dA(z')-\frac{c}{2\pi}\log|z|.
\end{equation}
By \eqref{eq:potential}, this function satisfies
\[
u\equiv 0
\qquad\text{on }\R^2\setminus (U\cup\{0\}).
\]

We first show that \(0\in U\). Since the singular term in
\eqref{eq:potato-second-proof} is \(-c\delta_0\), one must have \(0\in\overline U\). If
\(0\in\partial U\), then there exist exterior points \(z_j\to 0\) with
\[
u(z_j)=0.
\]
On the other hand, the first term in \eqref{eq:explicit-u-second-proof} remains finite as
\(z\to 0\), because \(f\in L^\infty(U)\) and \(\log|z-z'|\) is locally integrable, whereas the
second term equals \(-\frac{c}{2\pi}\log|z|\to +\infty\). This contradicts the exterior identity
\(u(z_j)=0\). Hence \(0\notin\partial U\), and therefore
\[
0\in U.
\]

We next show directly from \eqref{eq:explicit-u-second-proof} that
\[
u\in C^1(\R^2\setminus\{0\}).
\]
Indeed, if we write
\[
g(z):=f(|z|)\mathbf{1}_U(z)\in L^\infty_c(\R^2),
\]
then
\[
u(z)=\frac{1}{2\pi}\int_{\R^2}\log|z-z'|\,g(z')\,dA(z')-\frac{c}{2\pi}\log|z|.
\]
Since \(g\in L^\infty_c(\R^2)\), the only contribution to the integral defining the convolution
\((K\ast g)(z)\) below comes from a fixed bounded set, so we may replace \(K\) by its restriction
\(K\mathbf{1}_{B_R}\) for any sufficiently large ball \(B_R\). The kernel
\(K(\zeta)=\zeta/|\zeta|^2\) satisfies \(|K(\zeta)|=|\zeta|^{-1}\), which is integrable on bounded
subsets of \(\R^2\); hence \(K\mathbf{1}_{B_R}\in L^1(\R^2)\). It is a standard fact
that the convolution of an \(L^1\) function with a bounded function is continuous
(this follows from the density of \(C_c(\R^2)\) in \(L^1(\R^2)\) and continuity of translations in
\(L^1\)), so
\[
(K\ast g)(z):=\int_{\R^2}\frac{z-z'}{|z-z'|^2}\,g(z')\,dA(z')
\]
is a continuous function of \(z\in\R^2\). On the other hand, a direct distributional computation
shows
\[
\nabla\!\left(\int_{\R^2}\log|z-z'|\,g(z')\,dA(z')\right)=K\ast g
\qquad\text{in }\mathcal D'(\R^2),
\]
so the gradient of the convolution \(\log|\cdot|\ast g\) coincides almost everywhere with the
continuous function \(K\ast g\). It follows that \(\log|\cdot|\ast g\) is of class \(C^1\) on
\(\R^2\), and hence \(u\in C^1(\R^2\setminus\{0\})\) (the only obstruction to global \(C^1\)
regularity being the explicit logarithmic singularity at the origin coming from the
\(-\frac{c}{2\pi}\log|z|\) term).

Now let \(z_*\in\partial U\). Since \(0\in U\), the boundary point \(z_*\) is bounded away from
\(0\); fix a small ball \(B_\rho(z_*)\Subset \R^2\setminus\{0\}\). On the exterior open set
\(\R^2\setminus \overline U\) we have \(u\equiv 0\), so \(\nabla u\equiv 0\) there. As
\(\partial U\) is contained in the boundary of the exterior, every \(z_*\in\partial U\) is a limit
of exterior points \(z_j\to z_*\) along which \(\nabla u(z_j)=0\). By the \(C^1\) regularity of
\(u\) on \(B_\rho(z_*)\), the gradient \(\nabla u\) is continuous at \(z_*\), and we conclude
\(\nabla u(z_*)=0\). In complex notation \(\partial_z=\tfrac{1}{2}(\partial_x-i\partial_y)\), this
yields
\begin{equation}\label{eq:dz-zero-second-proof}
\partial_z u=0
\qquad\text{on }\partial U.
\end{equation}
We emphasize that \emph{no} boundary regularity of \(\partial U\) has been used at this point, in
contrast with the first proof: the vanishing of \(\nabla u\) on \(\partial U\) is forced solely by
exterior continuity, which in turn is supplied by the explicit
formula~\eqref{eq:explicit-u-second-proof}.

Define
\[
\Phi(t):=\int_0^t s^k e^{-\pi s}\,ds
\qquad\text{and}\qquad
\Psi(z):=\frac{\Phi(|z|^2)}{z},
\quad z\neq 0.
\]
A direct computation, using \(|z|^2=z\bar z\) and \(\partial_{\bar z}(z\bar z)=z\), gives, for
\(z\neq 0\),
\[
\partial_{\bar z}\Psi(z)
=
\frac{\Phi'(|z|^2)\,\partial_{\bar z}(|z|^2)}{z}
=
\frac{|z|^{2k}e^{-\pi|z|^2}\cdot z}{z}
=
|z|^{2k}e^{-\pi|z|^2}
=
f(|z|).
\]
At the origin, the apparent singularity \(1/z\) is integrable against the vanishing factor
\(\Phi(|z|^2)\): since \(k\ge 0\) and \(s^k e^{-\pi s}\le s^k\) for \(s\ge 0\), we have
\(\Phi(t)=\int_0^t s^k e^{-\pi s}\,ds\le \int_0^t s^k\,ds = t^{k+1}/(k+1)=O(t^{k+1})\) as
\(t\to 0\). Hence \(\Psi(z)=O(|z|^{2k+1})\) near \(z=0\), so \(\Psi\in L^\infty_{\mathrm{loc}}\)
(in fact continuous, with \(\Psi(0)=0\)). Consequently, the identity
\(\partial_{\bar z}\Psi=f(|z|)\) persists in the distributional sense on all of \(U\); no extra
\(\delta_0\) contribution arises at the origin.

Recall also the standard identity
\[
\partial_{\bar z}\!\left(\frac{1}{z}\right)=\pi\delta_0
\qquad\text{in }\mathcal D'(\R^2),
\]
which is equivalent to the fact that \(\frac{1}{\pi z}\) is a fundamental solution of
\(\partial_{\bar z}\); equivalently, \(\Delta\log|z|=4\partial_z\partial_{\bar z}\log|z|=2\pi\delta_0\)
combined with \(\partial_z\log|z|=\frac{1}{2z}\). Now set
\[
F(z):=4\partial_z u(z)-\Psi(z)+\frac{c}{\pi z},
\qquad z\in U\setminus\{0\}.
\]
We first check that \(F\) defines a distribution on all of \(U\) (a~priori it is given only on
\(U\setminus\{0\}\)). From \eqref{eq:explicit-u-second-proof} we have, near \(z=0\),
\[
4\partial_z u(z)
=
\frac{2}{\pi}\partial_z\!\left(\int_{\R^2}\log|z-z'|\,g(z')\,dA(z')\right)-\frac{c}{\pi z},
\]
where the first term is locally bounded near \(z=0\) (the Cauchy-type integral
\(\int (z-z')^{-1} g(z')\,dA(z')\) of a bounded compactly supported function is locally
integrable, by the same \(L^1_{\mathrm{loc}}\)-kernel argument as above). Adding \(-\Psi\)
(continuous with \(\Psi(0)=0\)) and \(c/(\pi z)\) cancels the explicit \(1/z\)-singularity, so
\(F\) is locally bounded near the origin and hence defines a distribution on \(U\).

Using \(\Delta=4\partial_z\partial_{\bar z}\), the relation \(\partial_{\bar z}(1/z)=\pi\delta_0\),
the computation of \(\partial_{\bar z}\Psi\) above, and \eqref{eq:potato-second-proof}, we
compute, in \(\mathcal D'(U)\),
\[
\partial_{\bar z}F
=
4\partial_z\partial_{\bar z}u-\partial_{\bar z}\Psi+\frac{c}{\pi}\partial_{\bar z}\!\left(\frac{1}{z}\right)
=
\Delta u-f(|z|)+c\,\delta_0
=
\bigl(f(|z|)-c\delta_0\bigr)-f(|z|)+c\delta_0
=0.
\]
The singular term \(-c\delta_0\) in \eqref{eq:potato-second-proof}, arriving through
\(4\partial_z\partial_{\bar z}u=\Delta u\), is thus exactly cancelled by
\((c/\pi)\partial_{\bar z}(1/z)=c\delta_0\). Hence \(F\) is a distributional solution of the
Cauchy--Riemann equation \(\partial_{\bar z}F=0\) on \(U\), and by Weyl's lemma for
\(\partial_{\bar z}\) (equivalently, hypoellipticity of the Cauchy--Riemann operator) it is
represented by a genuine holomorphic function on \(U\).

Define
\[
W(z):=\pi z\,F(z).
\]
Then \(W\) is holomorphic on \(U\) and \(W(0)=0\). Since \(\partial_z u\) is continuous up to
\(\partial U\), so is \(W\), and by \eqref{eq:dz-zero-second-proof},
\[
W|_{\partial U}
=
\pi z\left(4\partial_z u-\frac{\Phi(|z|^2)}{z}+\frac{c}{\pi z}\right)\Big|_{\partial U}
=
c-\pi\Phi(|z|^2).
\]
In particular, \(W|_{\partial U}\) is real-valued. The imaginary part of \(W\) is therefore
harmonic in \(U\), continuous up to \(\partial U\), and vanishes on \(\partial U\). By the
maximum principle,
\[
\Im W\equiv 0
\qquad\text{in }U.
\]
Since a holomorphic function with zero imaginary part is constant, there exists \(\kappa\in\R\)
such that
\[
W\equiv \kappa
\qquad\text{in }U.
\]

Restricting to the boundary, and using that \(W\) is continuous up to \(\partial U\) with
boundary value \(c-\pi\Phi(|z|^2)\), we obtain
\[
c-\pi\Phi(|z|^2)=\kappa
\qquad\text{for all }z\in\partial U.
\]
The function \(t\mapsto \Phi(t)=\int_0^t s^k e^{-\pi s}\,ds\) has derivative
\(\Phi'(t)=t^k e^{-\pi t}>0\) for \(t>0\), hence is strictly increasing on \([0,\infty)\). The
identity above therefore forces \(|z|^2\) to take a single value on \(\partial U\), i.e.\ \(|z|\)
is constant on \(\partial U\). Consequently \(\partial U\) is contained in a circle centered at
the origin. Let \(R>0\) be such that \(\partial U\subset \{|z|=R\}\). Since \(U\) is open and
\(U\cap\partial U=\varnothing\), we have \(U\subset \{|z|<R\}\cup\{|z|>R\}\); since \(U\) is
connected and the two pieces are disjoint open sets, \(U\) must lie entirely in one of them, and
boundedness of \(U\) forces \(U\subset \{|z|<R\}\). If this inclusion were strict, pick
\(z_1\in\{|z|<R\}\setminus U\) and \(z_0\in U\). The straight segment from \(z_0\) to \(z_1\) is
contained in the convex set \(\{|z|<R\}\), and as it joins a point of \(U\) to a point of its
complement it must cross \(\partial U\) at some \(z_*\); but then \(z_*\in\partial U\cap\{|z|<R\}\),
contradicting \(\partial U\subset \{|z|=R\}\). Therefore
\[
U=\{|z|<R\}
\]
that is, \(U\) is a disk.
\end{proof}

\begin{remark}
The two proofs use the boundary assumptions in rather different ways, and it is instructive to
keep careful track of the regularity bookkeeping.

In the first proof, once the logarithmic potential \(\Psi_U\) is written explicitly, no elliptic
boundary regularity is needed to conclude that \(\Psi_U\in C^1(\R^2\setminus\{0\})\): this follows
directly from convolution with the locally integrable kernel \(\nabla\log|z|\), exactly as in the
\(L^1_{\mathrm{loc}}\)-translation argument used here for \eqref{eq:explicit-u-second-proof}. Thus
the local radiality argument itself only uses continuity of the gradient up to the boundary. The
genuine extra input appears only in the final geometric step, where one needs that, near a suitable
point \(x_0\in \partial U\), the set \(\partial U\cap B_r(x_0)\) contains a nontrivial boundary
arc, so that the radii attained there form an interval. Accordingly, the first proof does not
really need \(C^2\) regularity; it works under any hypothesis guaranteeing such a local arc
structure, for instance when \(U\) is a Jordan domain. In particular, any Lipschitz, \(C^1\), or
\(C^2\) boundary is more than sufficient.

By contrast, the second proof does not use any boundary regularity at all. The explicit formula
\eqref{eq:explicit-u-second-proof} already yields \(u\in C^1(\R^2\setminus\{0\})\), and since
\(u\equiv 0\) on the exterior open set \(\R^2\setminus \overline U\), continuity implies that both
\(u\) and \(\nabla u\) vanish on \(\partial U\) -- precisely the boundary information packaged in
\eqref{eq:dz-zero-second-proof}. This is exactly the input needed to construct the holomorphic
function \(W=\pi z F\) and conclude that \(|z|\) is constant on \(\partial U\). Hence the second
proof recovers the theorem under the original assumptions of Abreu and D\"orfler
\cite{Abreu-Doerfler}, namely for arbitrary bounded simply connected domains, with no boundary
smoothness hypothesis. The mechanism is also conceptually closer to the classical
quadrature-domain rigidity results of \cite{Aharonov-Schiffer-Zalcman,Gustafsson,Karp}, where the
existence of a holomorphic function on \(U\) with real boundary values forces the boundary to lie
on a level set of a real-analytic function, here the elementary radial function
\(z\mapsto |z|^2\).
\end{remark}

% ==============================================================
% Section 5
% ==============================================================
